\chapter{Proportionnalité, pourcentages, échelles}

\noindent\textit{La proportionnalité est partout dans la vie courante : le prix payé au
marché selon la quantité achetée, une distance sur une carte, une vitesse constante, une
remise en pourcentage. Ce chapitre rassemble les outils pour reconnaître et exploiter une
situation de proportionnalité : produit en croix, pourcentages, évolutions successives,
échelles et vitesse moyenne.}
\vspace{8pt}

\connecteBanner

\section{Grandeurs proportionnelles}

\begin{definition}[Grandeurs proportionnelles]
Deux grandeurs sont \textbf{proportionnelles} si l'on passe des valeurs de l'une aux valeurs
de l'autre en multipliant toujours par le \textbf{même nombre}, appelé \textbf{coefficient de
proportionnalité}.
\end{definition}

\begin{propriete}[Produit en croix]
Dans un tableau de proportionnalité à deux colonnes de valeurs $a\,;b$ et $c\,;d$, on a
l'égalité des produits en croix :
\[
\frac ab=\frac cd \iff a\times d=b\times c.
\]
\end{propriete}

\begin{illus}
\begin{tikzpicture}
\draw[onbline,line width=0.8pt] (0,0) rectangle (6,1.6);
\draw[onbline,line width=0.8pt] (2,0) -- (2,1.6);
\draw[onbline,line width=0.8pt] (0,0.8) -- (6,0.8);
\node at (1,1.2) {\small Masse (kg)};
\node at (1,0.4) {\small Prix (FCFA)};
\node at (3,1.2) {\small $5$};
\node at (3,0.4) {\small $3\,500$};
\node at (5,1.2) {\small $8$};
\node at (5,0.4) {\small $x$};
\draw[onbo,line width=1.1pt] (3,1.2) -- (5,0.4);
\draw[onbo,line width=1.1pt] (5,1.2) -- (3,0.4);
\end{tikzpicture}
\fig{Figure — Produit en croix dans un tableau de proportionnalité :
$x=\dfrac{8\times3\,500}{5}=5\,600$ FCFA.}
\end{illus}

\begin{exemple}
Un sac de $5$ kg de riz coûte $3\,500$ FCFA. Le prix de $8$ kg vérifie
$\dfrac x8=\dfrac{3\,500}5$, donc $x=\dfrac{8\times3\,500}5=5\,600$ FCFA.
\end{exemple}

\section{Pourcentages}

\begin{definition}[Pourcentage]
Prendre $t\,\%$ d'une quantité $Q$, c'est calculer $\dfrac t{100}\times Q$.
\end{definition}

\begin{illus}
\begin{tikzpicture}
\draw[fill=onbsoft,draw=onbline,line width=0.8pt] (0,0) rectangle (10,1);
\draw[fill=onbo,draw=onbo,line width=0.8pt] (0,0) rectangle (1.8,1);
\draw[onbmuted,dashed] (1.8,-0.15) -- (1.8,1.3);
\node[above] at (1.8,1.35) {\small $18\,\%$};
\node[below] at (5,-0.3) {\small $45\,000$ FCFA $\ (100\,\%)$};
\end{tikzpicture}
\fig{Figure — $18\,\%$ de $45\,000$ FCFA représentés par la portion colorée :
$0{,}18\times45\,000=8\,100$ FCFA.}
\end{illus}

\begin{exemple}
$18\,\%$ de $45\,000$ FCFA $=0{,}18\times45\,000=8\,100$ FCFA.
\end{exemple}

\begin{propriete}[Taux de variation]
Le \textbf{taux de variation} entre une valeur initiale $V_i$ et une valeur finale $V_f$ est
\[
t=\frac{V_f-V_i}{V_i}\times100 \quad(\text{en }\%).
\]
Il est \textbf{positif} pour une augmentation, \textbf{négatif} pour une baisse.
\end{propriete}

\begin{exemple}
Un article passe de $25\,000$ à $29\,000$ FCFA : $t=\dfrac{29\,000-25\,000}{25\,000}\times100=16\,\%$
(hausse de $16\,\%$).
\end{exemple}

\section{Évolutions successives}

\begin{propriete}[Coefficient multiplicateur]
Augmenter une quantité de $t\,\%$ revient à la multiplier par $1+\dfrac t{100}$ ; la
diminuer de $t\,\%$ revient à la multiplier par $1-\dfrac t{100}$. Pour \textbf{deux}
évolutions successives, on \textbf{multiplie} les deux coefficients (on ne les additionne
\textbf{pas}).
\end{propriete}

\begin{exemple}
Un prix de $50\,000$ FCFA subit une hausse de $20\,\%$ puis une baisse de $15\,\%$ :
$50\,000\times1{,}20\times0{,}85=51\,000$ FCFA. Le coefficient global est
$1{,}20\times0{,}85=1{,}02$, soit une hausse globale de seulement $2\,\%$ (et non $5\,\%$).
\end{exemple}

\begin{remarque}
\textbf{Piège classique} : une hausse de $20\,\%$ suivie d'une baisse de $20\,\%$ ne ramène
\textbf{pas} au prix initial ! Coefficient global $1{,}20\times0{,}80=0{,}96$ : il reste une
baisse de $4\,\%$.
\end{remarque}

\section{Échelles}

\begin{definition}[Échelle]
L'\textbf{échelle} d'une carte ou d'un plan est le rapport (dans la \textbf{même unité}) entre
une distance sur le document et la distance réelle correspondante :
$\text{échelle}=\dfrac{\text{distance sur le document}}{\text{distance réelle}}$.
\end{definition}

\begin{illus}
\begin{tikzpicture}
\draw[fill=onbsoft,draw=onbo,line width=0.9pt] (0,0.8) rectangle (2,1.3);
\node[above] at (1,1.45) {\small carte : $4$ cm};
\draw[fill=onbo2,draw=onbo,line width=0.9pt] (0,0) rectangle (8,0.5);
\node[below] at (4,-0.2) {\small réalité : $1$ km};
\end{tikzpicture}
\fig{Figure — Sur une carte à l'échelle $1/25\,000$, $4$ cm représentent $1$ km dans la
réalité.}
\end{illus}

\begin{exemple}
Échelle $\dfrac1{25\,000}$ : $4$ cm sur la carte correspondent à
$4\times25\,000=100\,000$ cm $=1$ km dans la réalité.
\end{exemple}

\section{Vitesse moyenne}

\begin{definition}[Vitesse moyenne]
La \textbf{vitesse moyenne} est le quotient de la distance parcourue $d$ par la durée $t$
du trajet : $v=\dfrac dt$. On en déduit $d=v\times t$ et $t=\dfrac dv$.
\end{definition}

\begin{illus}
\begin{tikzpicture}
\begin{axis}[width=8cm,height=5.5cm,axis lines=middle,xlabel={\small temps (h)},ylabel={\small distance (km)},
  xmin=0,xmax=4,ymin=0,ymax=340,xtick={1,2,3},ytick={80,160,240,320}]
\addplot[onbo,line width=1.2pt,domain=0:3.5]{80*x};
\filldraw[onbink] (axis cs:3,240) circle(2pt);
\node[onbink,anchor=west] at (axis cs:3.1,240) {\small $(3\,;240)$};
\end{axis}
\end{tikzpicture}
\fig{Figure — À vitesse constante $80$ km/h, la distance parcourue est proportionnelle au
temps : $d=80t$.}
\end{illus}

\begin{exemple}
Une voiture parcourt $240$ km en $3$ h : $v=\dfrac{240}3=80$ km/h.
\end{exemple}

% ═══════════════════════════════════════════════════════════════
\section{L'essentiel à retenir}

\begin{synthese}
\textbf{Proportionnalité.} Coefficient constant ; produit en croix $ad=bc$.\\[4pt]
\textbf{Pourcentage.} $t\,\%$ de $Q=\dfrac t{100}\times Q$. Taux de variation :
$\dfrac{V_f-V_i}{V_i}\times100$.\\[4pt]
\textbf{Évolutions successives.} On \textbf{multiplie} les coefficients multiplicateurs
$\big(1\pm\frac t{100}\big)$ ; ne jamais additionner les pourcentages.\\[4pt]
\textbf{Échelle.} $\text{échelle}=\dfrac{\text{distance carte}}{\text{distance réelle}}$
(même unité).\\[4pt]
\textbf{Vitesse.} $v=\dfrac dt$ ;\ $d=v\times t$ ;\ $t=\dfrac dv$.\\[4pt]
\textbf{Piège fréquent.} $+20\,\%$ puis $-20\,\%$ ne redonne \textbf{pas} la valeur de
départ (coefficient global $1{,}2\times0{,}8=0{,}96\ne1$).
\end{synthese}

% ═══════════════════════════════════════════════════════════════
\section{Méthodes \& savoir-faire}

\methode{Compléter un tableau de proportionnalité}
Utiliser le produit en croix (ou calculer le coefficient de proportionnalité), puis
appliquer.
\begin{exemple}
Une recette pour $4$ personnes utilise $320$ g de farine. Pour $7$ personnes :
$\dfrac{320\times7}4=560$ g.
\end{exemple}

\methode{Calculer un pourcentage d'une quantité}
Multiplier la quantité par $\dfrac t{100}$.
\begin{exemple}
$35\,\%$ de $8\,000$ FCFA $=0{,}35\times8\,000=2\,800$ FCFA.
\end{exemple}

\methode{Calculer un taux de variation}
Calculer $\dfrac{V_f-V_i}{V_i}\times100$ ; le signe indique une hausse ($+$) ou une baisse
($-$).
\begin{exemple}
Une population passe de $40\,000$ à $34\,000$ habitants :
$\dfrac{34\,000-40\,000}{40\,000}\times100=-15\,\%$ (baisse de $15\,\%$).
\end{exemple}

\methode{Appliquer une évolution en pourcentage}
Multiplier par $1+\dfrac t{100}$ (hausse) ou $1-\dfrac t{100}$ (baisse).
\begin{exemple}
Un salaire de $120\,000$ FCFA augmente de $8\,\%$ : nouveau salaire
$=120\,000\times1{,}08=129\,600$ FCFA.
\end{exemple}

\methode{Enchaîner deux évolutions successives}
Multiplier les deux coefficients multiplicateurs (jamais additionner les pourcentages).
\begin{exemple}
Une hausse de $10\,\%$ puis une baisse de $5\,\%$ : coefficient global
$1{,}10\times0{,}95=1{,}045$, soit une hausse globale de $4{,}5\,\%$.
\end{exemple}

\methode{Utiliser une échelle}
Passer de la distance sur le document à la distance réelle (ou l'inverse) en multipliant ou
divisant par l'échelle, en veillant aux unités.
\begin{exemple}
Distance réelle $15$ km sur une carte à l'échelle $\dfrac1{200\,000}$ : distance sur la
carte $=\dfrac{15\,\text{km}}{200\,000}=\dfrac{1\,500\,000\text{ cm}}{200\,000}=7{,}5$ cm.
\end{exemple}

\methode{Calculer une vitesse moyenne, une distance ou une durée}
Utiliser $v=\dfrac dt$, $d=v\times t$ ou $t=\dfrac dv$ selon la grandeur cherchée, en gardant
des unités cohérentes.
\begin{exemple}
Un cycliste roule à $18$ km/h pendant $2$ h $30$ (soit $2{,}5$ h) :
$d=18\times2{,}5=45$ km.
\end{exemple}

% ═══════════════════════════════════════════════════════════════
\section{Exercices d'application}
\begin{multicols}{2}

\rubrique{Tableaux de proportionnalité}
\exo{}\diff{1} Un sac de $5$ kg de riz coûte $3\,500$ FCFA. Calculer le prix de $8$ kg.

\exo{}\diff{1} Une recette pour $4$ personnes utilise $320$ g de farine. Quelle quantité
faut-il pour $7$ personnes ?

\exo{}\diff{2} $12$ stylos coûtent $3\,000$ FCFA. Calculer le prix de $20$ stylos.

\exo{}\diff{2} Une voiture consomme $42$ L d'essence pour $600$ km. Quelle sera sa
consommation pour $850$ km (arrondir au dixième de litre) ?

\rubrique{Pourcentages}
\exo{}\diff{1} Calculer $18\,\%$ de $45\,000$ FCFA.

\exo{}\diff{1} Calculer $35\,\%$ de $8\,000$ FCFA.

\exo{}\diff{2} Un article passe de $25\,000$ FCFA à $29\,000$ FCFA. Calculer le taux de
variation.

\exo{}\diff{2} Une population passe de $40\,000$ à $34\,000$ habitants. Calculer le taux de
variation.

\rubrique{Évolutions successives}
\exo{}\diff{2} Un prix de $50\,000$ FCFA subit une hausse de $20\,\%$ puis une baisse de
$15\,\%$. Calculer le prix final, et le taux d'évolution global.

\exo{}\diff{2} Calculer le coefficient multiplicateur global correspondant à une hausse de
$20\,\%$ suivie d'une baisse de $15\,\%$.

\exo{}\diff{1} Un salaire de $120\,000$ FCFA augmente de $8\,\%$. Calculer le nouveau
salaire.

\exo{}\diff{3} Un article soldé à $30\,\%$ coûte maintenant $21\,000$ FCFA. Quel était son
prix initial ?

\rubrique{Échelles et vitesse}
\exo{}\diff{1} Sur une carte à l'échelle $\dfrac1{25\,000}$, une distance mesure $4$ cm.
Calculer la distance réelle, en km.

\exo{}\diff{2} Une distance réelle de $15$ km est représentée par quelle longueur sur une
carte à l'échelle $\dfrac1{200\,000}$ ?

\exo{}\diff{1} Une voiture parcourt $240$ km en $3$ h. Calculer sa vitesse moyenne.

\exo{}\diff{2} Un cycliste roule à $18$ km/h pendant $2$ h $30$. Quelle distance
parcourt-il ?

% ═══════════════════════════════════════════════════════════════
\end{multicols}

\section{Exercices d'entraînement}
\begin{multicols}{2}

\rubrique{Tableaux de proportionnalité}
\exo{}\diff{1} Un sac de $2$ kg de sucre coûte $250$ FCFA. Calculer le prix de $5$ kg.

\exo{}\diff{2} Une recette pour $6$ personnes utilise $450$ g de sucre. Quelle quantité
faut-il pour $9$ personnes ?

\exo{}\diff{2} $15$ cahiers coûtent $3\,000$ FCFA. Calculer le prix de $36$ cahiers.

\exo{}\diff{2} Une moto parcourt $270$ km avec $8$ L d'essence. Quelle distance
parcourt-elle avec $3$ L ?

\rubrique{Pourcentages}
\exo{}\diff{1} Calculer $12\,\%$ de $60\,000$ FCFA.

\exo{}\diff{1} Calculer $40\,\%$ de $9\,500$ FCFA.

\exo{}\diff{2} Un loyer passe de $18\,000$ à $21\,600$ FCFA. Calculer le taux de variation.

\exo{}\diff{2} Un troupeau passe de $50\,000$ à $42\,000$ têtes de bétail. Calculer le
taux de variation.

\rubrique{Évolutions successives}
\exo{}\diff{2} Un salaire de $80\,000$ FCFA augmente de $10\,\%$ une année, puis de
$5\,\%$ l'année suivante. Calculer le salaire final et le taux d'évolution global.

\exo{}\diff{2} Calculer le coefficient multiplicateur global correspondant à deux baisses
successives de $10\,\%$.

\exo{}\diff{1} Un prix de $200\,000$ FCFA augmente de $5\,\%$. Calculer le nouveau prix.

\exo{}\diff{3} Après une baisse de $13\,\%$, un article coûte $17\,400$ FCFA. Quel était
son prix initial ?

\rubrique{Échelles}
\exo{}\diff{2} Sur un plan à l'échelle $\dfrac1{50\,000}$, une distance mesure $6$ cm.
Calculer la distance réelle, en km.

\exo{}\diff{2} Une distance réelle de $20$ km est représentée par quelle longueur sur une
carte à l'échelle $\dfrac1{500\,000}$ ?

\exo{}\diff{2} Sur un plan à l'échelle $\dfrac1{200}$, une pièce mesure $6$ cm. Calculer sa
longueur réelle en mètres.

\rubrique{Vitesse}
\exo{}\diff{1} Un bus parcourt $315$ km en $4$ h $30$. Calculer sa vitesse moyenne.

\exo{}\diff{2} Un piéton marche à $4{,}5$ km/h pendant $50$ min. Quelle distance
parcourt-il (arrondir au centième de km) ?

\exo{}\diff{2} Un camion roule à $60$ km/h. Combien de temps met-il pour parcourir $150$
km ?

\exo{}\diff{2} Un avion parcourt $2\,100$ km en $3$ h. Calculer sa vitesse moyenne, puis le
temps nécessaire pour parcourir $3\,500$ km à cette même vitesse.

% ═══════════════════════════════════════════════════════════════
\end{multicols}

\section{Sujets type BEPC}

\sujetbac[proportionnalité et pourcentages au marché — Bafoussam]
\begin{enumerate}[label=\arabic*)]
\item Un sac de $25$ kg de maïs coûte $8\,750$ FCFA. Calculer le prix de $40$ kg de maïs
(en utilisant un produit en croix).
\item Un commerçant de Bafoussam augmente ses prix de $12\,\%$. Un article coûtait
$15\,000$ FCFA. Calculer son nouveau prix.
\item Le prix d'un autre article est passé de $20\,000$ FCFA à $23\,000$ FCFA. Calculer le
taux de variation.
\item Après une hausse de $10\,\%$ suivie d'une baisse de $5\,\%$, le prix d'un sac de riz
est de $41\,800$ FCFA.
\begin{enumerate}[label=\alph*)]
\item Calculer le coefficient multiplicateur global de ces deux évolutions.
\item En déduire le prix initial du sac de riz.
\end{enumerate}
\end{enumerate}

\sujetbac[échelle, vitesse et remise]
\begin{enumerate}[label=\arabic*)]
\item Sur une carte à l'échelle $\dfrac1{50\,000}$, la distance entre deux villages mesure
$20$ cm. Calculer la distance réelle, en km.
\item Un bus parcourt cette distance en $15$ minutes. Calculer sa vitesse moyenne en km/h.
\item Un article vendu $8\,000$ FCFA est soldé à $6\,000$ FCFA. Calculer le taux de
réduction appliqué.
\item Une population de $24\,000$ habitants augmente de $5\,\%$ par an.
\begin{enumerate}[label=\alph*)]
\item Calculer la population après $1$ an.
\item Calculer la population après $2$ ans (arrondir à l'unité).
\end{enumerate}
\end{enumerate}

\sujetbac[pourcentages, évolutions et échelle — vie scolaire]
\begin{enumerate}[label=\arabic*)]
\item Une classe de $45$ élèves compte $60\,\%$ de filles. Calculer le nombre de filles et
le nombre de garçons.
\item Un club sportif compte $250$ membres. Son effectif augmente de $8\,\%$ une année,
puis de $5\,\%$ l'année suivante.
\begin{enumerate}[label=\alph*)]
\item Calculer le coefficient multiplicateur global de ces deux évolutions.
\item Calculer l'effectif du club après ces deux augmentations (arrondir à l'unité).
\end{enumerate}
\item Le prix d'une inscription est passé de $12\,000$ FCFA à $13\,800$ FCFA. Calculer le
taux d'augmentation.
\item Sur un plan à l'échelle $\dfrac1{200}$, un couloir mesure $8$ cm. Calculer sa
longueur réelle, en mètres.
\end{enumerate}

% ═══════════════════════════════════════════════════════════════
\section{Corrigés}
\begin{multicols}{2}

\corrige{de l'exercice 1}
$\dfrac{8\times3\,500}5=5\,600$ FCFA.

\corrige{de l'exercice 2}
$\dfrac{320\times7}4=560$ g.

\corrige{de l'exercice 3}
$\dfrac{3\,000\times20}{12}=5\,000$ FCFA.

\corrige{de l'exercice 4}
$\dfrac{42\times850}{600}=59{,}5$ L.

\corrige{de l'exercice 5}
$0{,}18\times45\,000=8\,100$ FCFA.

\corrige{de l'exercice 6}
$0{,}35\times8\,000=2\,800$ FCFA.

\corrige{de l'exercice 7}
$\dfrac{29\,000-25\,000}{25\,000}\times100=16\,\%$.

\corrige{de l'exercice 8}
$\dfrac{34\,000-40\,000}{40\,000}\times100=-15\,\%$.

\corrige{de l'exercice 9}
$50\,000\times1{,}20\times0{,}85=51\,000$ FCFA. Taux global :
$\dfrac{51\,000-50\,000}{50\,000}\times100=2\,\%$.

\corrige{de l'exercice 10}
$1{,}20\times0{,}85=1{,}02$.

\corrige{de l'exercice 11}
$120\,000\times1{,}08=129\,600$ FCFA.

\corrige{de l'exercice 12}
Prix initial $\times0{,}70=21\,000$, donc prix initial $=\dfrac{21\,000}{0{,}70}=30\,000$
FCFA.

\corrige{de l'exercice 13}
$4\times25\,000=100\,000$ cm $=1$ km.

\corrige{de l'exercice 14}
$15$ km $=1\,500\,000$ cm ; $\dfrac{1\,500\,000}{200\,000}=7{,}5$ cm.

\corrige{de l'exercice 15}
$v=\dfrac{240}3=80$ km/h.

\corrige{de l'exercice 16}
$d=18\times2{,}5=45$ km.

\corrige{du sujet 1}
1) $\dfrac{8\,750\times40}{25}=14\,000$ FCFA.\\
2) $15\,000\times1{,}12=16\,800$ FCFA.\\
3) $\dfrac{23\,000-20\,000}{20\,000}\times100=15\,\%$.\\
4a) $1{,}10\times0{,}95=1{,}045$.\\
4b) Prix initial $=\dfrac{41\,800}{1{,}045}=40\,000$ FCFA.

\corrige{du sujet 2}
1) $20\times50\,000=1\,000\,000$ cm $=10$ km.\\
2) $15$ min $=0{,}25$ h ; $v=\dfrac{10}{0{,}25}=40$ km/h.\\
3) $\dfrac{8\,000-6\,000}{8\,000}\times100=25\,\%$.\\
4a) $24\,000\times1{,}05=25\,200$ habitants.\\
4b) $24\,000\times1{,}05^{2}=26\,460$ habitants.

\corrige{du sujet 3}
1) Filles : $45\times0{,}60=27$ ; garçons : $45-27=18$.\\
2a) $1{,}08\times1{,}05=1{,}134$.\\
2b) $250\times1{,}134=283{,}5\approx284$ membres.\\
3) $\dfrac{13\,800-12\,000}{12\,000}\times100=15\,\%$.\\
4) $8\times200=1\,600$ cm $=16$ m.

\end{multicols}
\exercicesBanner
